\documentclass[10pt,a4paper]{amsart}
\usepackage[margin=19mm]{geometry}
\usepackage{amsmath,amssymb,mathtools}
\usepackage[scr=rsfs,cal=euler]{mathalfa}
\usepackage{xfrac}
\usepackage[unicode,hidelinks,bookmarks=false]{hyperref}
\newcommand{\ie}{i.e.\ }
\newcommand{\cf}{cf.\ }
\newcommand{\resp}{respectively\ }
\newcommand{\Subsection}{\subsection}
\providecommand{\nobreakdash}{\nobreak\hbox{-}}
\setlength{\emergencystretch}{2em}
\multlinegap0pt
\usepackage{bm}
\usepackage{bbm}
\usepackage{dsfont}
\usepackage{xspace}
\usepackage{cancel}
\usepackage{mathtools}
\usepackage[shortlabels]{enumitem}
\usepackage{amsthm}
\makeatletter
\@ifundefined{theorem}{\newtheorem{theorem}{Theorem}}{}
\@ifundefined{lemma}{\newtheorem{lemma}[theorem]{Lemma}}{}
\@ifundefined{proposition}{\newtheorem{proposition}[theorem]{Proposition}}{}
\makeatother
\makeatletter
\newcommand{\R}{\mathbb{R}}
\newcommand{\cX}{\mathcal{X}}
\newcommand{\ep}{\epsilon}
\expandafter\ifx\csname proof\endcsname\relax
\newenvironment{proof}[1][Proof]{\noindent\textbf{#1.} }{\ \rule{0.5em}{0.5em}}
\fi
\newcommand{\red}[1]{{#1}}
\makeatother
\title[Scaling Limit of the Kuramoto Model on Random Geometric Grap]{Scaling Limit of the Kuramoto Model on Random Geometric Graphs}
\author{Francisco Cirelli, Pablo Groisman, Ruojun Huang,, Hern\'an Vivas}
\date{Source: arXiv:2402.15311v3 (CC BY 4.0)}
\begin{document}
\maketitle

\noindent\emph{Source and licence.} Scaling Limit of the Kuramoto Model on Random Geometric Graphs. Version arXiv:2402.15311v3 (CC BY 4.0), distributed under the Creative Commons Attribution 4.0 International licence, as recorded in the arXiv licence metadata for this exact version and shown on the abstract page https://arxiv.org/abs/2402.15311v3. Full rights notice: licenses/arxiv-2402.15311-CC-BY-4.0.txt.\par\smallskip

\noindent\textit{Editorial scope.} Counted unit: the Proposition whose source label is \texttt{prop.existence} in the article (source0.tex lines 614--616, source label \texttt{prop.existence}), delivered together with the article's own complete original proof of it (source0.tex lines 618--704, one \texttt{proof} environment of 87 lines). No mathematics is written, repaired or re-derived: statement and proof are the article's own text line for line. Required context retained uncounted: def:mod (lines 333--335); eq.integral (lines 530--535); eq.fixedpoint (lines 549--551); lem:derivative (lines 558--564); lem:periodic (lines 595--601). Every definition, display and label of those lines is retained unchanged. Omitted from this package and not redistributed: 1--332, 336--529, 536--548, 552--557, 565--594, 602--613, 617--617, 705--1340. Nothing inside a retained excerpt is omitted or altered: every retained body is delivered exactly as the article's lines are written, so no omission receipt of any kind accompanies this package. Citations: the retained excerpts carry no citation command at all, so no bibliography entry of the article is transcribed and no reference is redistributed. Reference adaptation: none was needed and none was made; every reference inside the delivered unit resolves inside it, so no unresolved reference or citation is produced. Printed slips kept literal: no printed word, symbol, hypothesis or number of the counted statement or of its proof was altered. Layout and class: the article's own class is \texttt{amsart} with packages including amsmath, amsthm, amssymb, amsfonts, caption, enumitem, soul, hyperref, color, inputenc, babel, tikz, url, natbib; the wrapper is the lane's pinned \texttt{amsart} editorial wrapper at 10pt on A4, which restates the theorem-like environments the retained text needs and adds the article's own macro definitions that the retained text needs (\texttt{\textbackslash R}, \texttt{\textbackslash cX}, \texttt{\textbackslash ep}, \texttt{\textbackslash red}), with extra wrapper packages (bm, bbm, dsfont, xspace, cancel, mathtools, enumitem). The delivered numbering is the unit's own. Assets: no retained span contains a figure environment, a graphics-inclusion command, a PDF-page inclusion, a table or a TikZ picture; no image file is copied into this package, the package declares no asset dependency and no image file is redistributed with it. Licence: version arXiv:2402.15311v3 of this article is distributed under the Creative Commons Attribution 4.0 International licence, as recorded in the arXiv licence metadata for this exact version and shown on the abstract page https://arxiv.org/abs/2402.15311v3. A full rights notice is delivered at licenses/arxiv-2402.15311-CC-BY-4.0.txt. Characters: every retained excerpt is pure ASCII, so the delivered engine, XeLaTeX with its default Latin Modern text font, reports no missing character.
\begin{align}\label{def:mod}
[y]:=y \text{ mod }2\pi \quad\in\Lambda_d,
\end{align}

\begin{eqnarray}\label{eq.integral}
\begin{cases}
\displaystyle{\frac{d}{dt}}u^{I,\ep}(t,x)=\frac{1}{\sigma_d\ep^{d+2}}\int_{\red{\R^d}}\sin\left(u^{I,\ep}(t,y)-u^{I,\ep}(t,x)\right)K\Big(\frac{y-x}{\ep}\Big)dy, \label{eq:int}\\[10pt]
u^{I,\ep}|_{t=0}=\red{\tilde u_0,}
\end{cases}
\end{eqnarray}

\begin{equation}\label{eq.fixedpoint}
u^{I,\ep}(t,x) =\tilde u_0(x) + \frac{1}{\sigma_d\ep^{d+2}}\int_0^t\int_{\red{\R^d}}\sin\left(u^{I,\ep}(s,y) - u^{I,\ep}(s,x)\right)K\Big(\frac{y-x}{\ep}\Big)dy\:ds.
\end{equation}

\begin{lemma}\label{lem:derivative}
Let $F(x):=\int_{\red{\R^d}}g(x,y)\mathsf K(y-x)dy$, where 
$g(x,y)\in C^1_b(\R^d\times\R^d,\R)$ \red{(with bounded $C^1$-norm)} and $\mathsf K\in L^1(\R^d,\R)$ with compact support, then the following derivative formula holds:
\begin{align*}
\nabla F(x)=\int_{\red{\R^d}}\left[\nabla_x g(x,y)+\nabla_yg(x,y) \right] \mathsf K(y-x) dy.
\end{align*}
\end{lemma}

\begin{lemma}\label{lem:periodic}
Suppose that $u(t,x):[0,\infty)\times\R^d\to\R$ is a continuous function in $(t,x)$, pseudo-periodic in $x$ for every $t$, then the function
\[
U(t,x):=\int_{\R^d}\sin\left(u(t,y) - u(t,x)\right)K\Big(\frac{y-x}{\ep}\Big)dy
\]
is $2\pi$-periodic in $x$ in each coordinate direction, for every $t$.
\end{lemma}

\begin{proposition}\label{prop.existence}
Fix any $\ep>0$. Given \red{any pseudo-periodic $\tilde u_0\in C^1(\R^d,\R)$}, there exists a unique \red{pseudo-periodic function $u^{I,\ep}\in C\left([0,\infty);C^1(\R^d,\R)\right)$} satisfying \eqref{eq.fixedpoint} and hence a unique solution of \eqref{eq.integral}.
\end{proposition}

\begin{proof}

Solutions of \eqref{eq.fixedpoint} are fixed points of the operator
\begin{align}\label{fixed-pt-map}
F_{\tilde u_0}(u^{I,\ep})(t,x) :=\tilde  u_0(x) + \frac{1}{\sigma_d\ep^{d+2}}\int_0^t\int_{\red{\R^d}}\sin\left(u^{I,\ep}(s,y) - u^{I,\ep}(s,x)\right)K\Big(\frac{y-x}{\ep}\Big)dy\:ds.
\end{align}
For a fixed \red{pseudo-periodic} initial condition \red{$\tilde u_0\in C^1(\R^d,\R)$} and a positive $T$ (to be chosen later) we consider the Banach space 
\red{\begin{align}\label{banach-ball}
\mathcal{X}_T:=\Big\{f: [0,T]\times\R^d\to\R:\; & f(t,\cdot)\text{ is pseudo-periodic},\, \, f\in C\big([0,T];C^1(\R^d,\R)\big), \nonumber\\
&f|_{t=0}=\tilde u_0, \;  \sup_{t\in[0,T]}\|f(t,\cdot)\|_{C^1(\Lambda_d)}\le 1+\|\tilde u_0\|_{C^1(\Lambda_d)}\Big\}
\end{align}
with the norm 
\[
\|f\|_{\mathcal{X}_T} := \sup_{t\in[0,T]}\|f(t,\cdot)\|_{C^1(\Lambda_d)},
\]
where by an abuse of notation, for a pseudo-periodic function $g:\R^d\to\R$, we write $\|g\|_{C^1(\Lambda_d)}$ for $\|g|_{\Lambda_d}\|_{C^1(\Lambda_d)}$. We immediately remark that the map $F_{\tilde u_0}$ \eqref{fixed-pt-map} takes a pseudo-periodic continuous function to another pseudo-periodic continuous function, with the same winding numbers. This is a consequence of Lemma \ref{lem:periodic} and that $\tilde u_0$ is pseudo-periodic.
}

We want to apply Banach's fixed point theorem  to $F_{\tilde u_0}$ in $\mathcal{X}_T$. We must check:

\begin{enumerate}[label=(\alph*)]

\item $F_{\tilde u_0}(\mathcal{X}_T)\subset\mathcal{X}_T$;

\item $F_{\tilde u_0}$ is a contraction, i.e.~there exists $\nu\in(0,1)$ such that 
\begin{equation*}
\|F_{\tilde u_0}(u^{I,\ep})-F_{\tilde u_0}(v)\|_{\mathcal{X}_T}\leq \nu\|u-v\|_{\mathcal{X}_T}
\end{equation*}
for all $u, v\in \mathcal{X}_T$.
\end{enumerate}

We first show (b). \red{We first observe that $h(x,y):=\sin\left(u(t,y)-u(t,x)\right)$ has bounded $C^1(\R^{2d})$-norm for each $t$. Indeed, $\|h\|_{L^\infty(\R^{2d})}\le 1$ and furthermore, $$\nabla_yh(x,y)=\cos\left(u(t,y)-u(t,x)\right)\nabla_y u(t,y).$$ Since $u(t,y)$ is pseudo-periodic, $\|\nabla_y u(t,y)\|_{L^\infty(\R^d)}$ is equal to $\|\nabla_y u(t,y)\|_{L^\infty(\Lambda_d)}$, and therefore is bounded. Hence $\|\nabla_yh(x,y)\|_{L^\infty(\R^{2d})}$ is bounded, and similarly $\|\nabla_xh(x,y)\|_{L^\infty(\R^{2d})}$ is bounded as well.}
We are in a situation to apply Lemma \ref{lem:derivative}, whereby for any $x\in\Lambda_d$,
%\begin{footnotesize}
\begin{align*}
& \sigma_d\ep^{d+2}|\nabla F_{\tilde u_0}(u)(t,x)-\nabla F_{\tilde u_0}(v)(t,x)| \\
& \le \int_0^t\int_{\R^d}\Big|\cos\left(u(s,y)-u(s,x)\right)\nabla u(s,y)-\cos\left(v(s,y)-v(s,x)\right)\nabla v(s,y)\Big|K\Big(\frac{y-x}{\ep}\Big)dy\:ds  \\
&+ \int_0^t\int_{\R^d}\Big|\cos\left(u(s,y)-u(s,x)\right)\nabla u(s,x)-\cos\left(v(s,y)-v(s,x)\right)\nabla v(s,x)\Big|K\Big(\frac{y-x}{\ep}\Big)dy\:ds.
\end{align*}
\red{
We notice that although we consider $x\in\Lambda_d$, there may be some $y\in\R^d$ that are within Euclidean distance $\ep$ from $x$ but are outside of $\Lambda_d$. Since $\ep\in(0,1)$, such $y$ must be in one of the adjacent boxes. We must have either $y_\ell=[y]_\ell\pm 2\pi$ or $y_\ell=[y]_\ell$, $\ell=1,2,...,d$ (see \eqref{def:mod}). For $g$ that is pseudo-periodic with winding numbers $\{k_\ell\}_{\ell=1}^d$, we have $g([y])=g(y)+\sum_{\ell=1}^da_\ell k_\ell \vec e_\ell$ for some $a_\ell\in\{0,\pm 2\pi\}$ and $\nabla g([y])=\nabla g(y)$. Hence, since the cosine is a $2\pi$-periodic function and $u,v$ are pseudo-periodic, the previous display equals
\begin{align*}
    %& \sigma_d\ep^{d+2}|\nabla F_{\tilde u_0}(u)(t,x)-\nabla F_{\tilde u_0}(v)(t,x)| \\
& =\int_0^t\int_{\R^d}\Big|\cos\left(u(s,[y])-u(s,x)\right)\nabla u(s,[y])-\cos\left(v(s,[y])-v(s,x)\right)\nabla v(s,[y])\Big|K\Big(\frac{y-x}{\ep}\Big)dy\:ds  \\
&+ \int_0^t\int_{\R^d}\Big|\cos\left(u(s,[y])-u(s,x)\right)\nabla u(s,x)-\cos\left(v(s,[y])-v(s,x)\right)\nabla v(s,x)\Big|K\Big(\frac{y-x}{\ep}\Big)dy\:ds.
\end{align*}}

Then by the triangle inequality and $1$-Lipschitz property of sine function, we further bound it by
\begin{align*}
  % & \sigma_d\ep^{d+2}|\nabla F_{\tilde u_0}(u)(t,x)-\nabla F_{\tilde u_0}(v)(t,x)| \\
& \le \int_0^t\int_{\R^d}\Big(|\nabla u(s,[y])-\nabla v(s,[y])| \\
&\hspace{4cm}+\big(|u(s,[y])-v(s,[y])|+|u(s,x)-v(s,x)|\big)|\nabla v(s,[y])|\Big)K\Big(\frac{y-x}{\ep}\Big)dy\:ds\\
&+\int_0^t\int_{\R^d}\Big(|\nabla u(s,x)-\nabla v(s,x)| \\
&\hspace{4cm}+\big(|u(s,[y])-v(s,[y])|+|u(s,x)-v(s,x)|\big)|\nabla v(s,x)|\Big)K\Big(\frac{y-x}{\ep}\Big)dy\:ds.
\end{align*}


Since $\|K\|_{L^1(\R^d)}=1$, applying \red{H\"older's inequality $$\|fK(\frac{y - \cdot}{\ep})\|_{L^1(\R^d)}\le \|f\|_{L^\infty(\R^d)}\|K(\frac{y-\cdot}{\ep})\|_{L^1(\R^d)} = \ep^d\|f\|_{L^\infty(\R^d)},$$ to the above integrals in $dy$, and noting that both $x,[y]\in\Lambda_d$}, we get from above
\begin{align}\label{eq:grad-diff}
&\sup_{t\in[0,T]}\|\nabla F_{\tilde u_0}(u)(t,\cdot)-\nabla F_{\tilde u_0}(v)(t,\cdot)\|_{L^\infty(\Lambda_d)}\\
&\le C\left(1+\|\tilde u_0\|_{C^1(\Lambda_d)}\right)T \sup_{s\in[0,T]}\|u(s,\cdot)-v(s,\cdot)\|_{L^\infty(\Lambda_d)}+CT \sup_{s\in[0,T]}\|\nabla u(s,\cdot)-\nabla v(s,\cdot)\|_{L^\infty(\Lambda_d)}, \nonumber
\end{align} 
where the constant $C=C(\ep, d)$ may change from line to line, and we used that since $v\in \cX_T$ \eqref{banach-ball}, its $C^1(\Lambda_d)$-norm in space is bounded by $1+\|\tilde u_0\|_{C^1(\Lambda_d)}$ up to time $T$.

Similarly, we can estimate 
\begin{align*}
&\sigma_d\ep^{d+2} |F_{\tilde u_0}(u)(t,x)-F_{\tilde u_0}(v)(t,x)| \\
& \leq \int_0^t\int_{\R^d}\Big|\sin\left(u(s,y)-u(s,x)\right)-\sin\left(v(s,y)-v(s,x)\right)\Big|K\Big(\frac{y-x}{\ep}\Big)dy\:ds \\
&= \red{\int_0^t\int_{\R^d}\Big|\sin\left(u(s,[y])-u(s,x)\right)-\sin\left(v(s,[y])-v(s,x)\right)\Big|K\Big(\frac{y-x}{\ep}\Big)dy\:ds,}
\end{align*}
which yields that
\begin{align}\label{eq:diff}
&\sup_{t\in[0,T]}\|F_{\tilde u_0}(u)(t,\cdot)-F_{\tilde u_0}(v)(t,\cdot)\|_{L^\infty(\Lambda_d)}\le C(\ep, d)T \sup_{t\in[0,T]}\|u(s,\cdot)-v(s,\cdot)\|_{L^\infty(\Lambda_d)}.
\end{align}
Recalling the definition of $C^1(\Lambda_d)$-norm, we get from \eqref{eq:grad-diff}-\eqref{eq:diff}
\begin{align*}
\|F_{\tilde u_0}(u)-F_{\tilde u_0}(v)\|_{\cX_T}\le CT\left(1+\|\tilde u_0\|_{C^1(\Lambda_d)}\right)\|u-v\|_{\cX_T},
\end{align*}
for some finite $C=C(\ep, d)$. 

For requirement (a), a similar argument (just do not consider $v$) shows that 
\begin{align*}
\|F_{\tilde u_0}(u)\|_{\cX_T}\le \|\tilde u_0\|_{C^1(\Lambda_d)}+CT\|u\|_{\cX_T}.
\end{align*}

Choosing $T_0:=\frac{1}{2C\left(1+\|\tilde u_0\|_{C^1(\Lambda_d)}\right)}$ we get by Banach's fixed point theorem local existence of a unique solution to \eqref{eq.fixedpoint} in the time interval $[0,T_0]$. We want to iterate the above argument and get global existence. We note that even though $T_0$ depends inversely on $C^1$-norm of the initial condition, the dependence is linear. On the other hand, our construction \eqref{banach-ball} is such that each iteration only increases the norm of the solution by at most $1$. This means that, if we iterate the previous argument infinitely many times, the sum of the local existence intervals diverges and we can reach any $T$. Finally note that a continuous solution of \eqref{eq.fixedpoint} is in fact differentiable in $t$ and as such, a solution of \eqref{eq.integral}.
\end{proof}

\end{document}
